% update_ch2_nof.tex — Описание NOF эксперимента и данных
% Каноническое описание экспериментальных данных кальциевого имиджинга.
% Используется в главах 2 и 3. Подключается через \input в начале главы 2.

\section{Описание эксперимента и данных}\label{sec:nof_data}

\subsection{Животные}

В экспериментах использовались мыши линии C57BL/6 обоих полов ($n=16$), в возрасте 3--5 месяцев. До хирургического вмешательства животные содержались в стандартных лабораторных клетках (по 2--7 особей) при 12-часовом цикле свет/темнота со свободным доступом к пище и воде. После операции мыши содержались индивидуально, прочие условия содержания оставались неизменными. Все эксперименты проводились в светлую фазу суток (с 10:00 до 18:00).

\subsection{Хирургические процедуры и доставка вирусного вектора}

Все процедуры были одобрены Комиссией по биоэтике Московского государственного университета имени М.\,В.~Ломоносова (заявка № 159-g, утверждённая на заседании Комиссии по биоэтике № 154-d-r от 17.08.2023) и проводились в соответствии с приказом МЗ РФ № 267 и Руководством NIH по содержанию и использованию лабораторных животных.

Хирургические процедуры выполнялись в три этапа: инъекция вирусного вектора, имплантация GRIN-линзы и установка базовой пластины. Наркоз индуцировался внутрибрюшинной инъекцией золетила (40~мг/кг) и ксилазина (5~мг/кг). За 5 минут до начала операции подкожно вводился дексаметазон (4~мг/кг) для уменьшения воспаления и предотвращения отёка мозга. Для предотвращения высыхания глаз наносился увлажняющий гель. Мыши фиксировались в стереотаксической раме (Kopf, США), температура тела поддерживалась с помощью обогревающей подкладки (Physitemp, США).

Для инъекции вируса 300~нл раствора AAV1.CAG.GCaMP6s вводилось унилатерально в область дорсального CA1 (координаты: AP $-1{,}94$~мм, ML $+1{,}46$~мм, DV $-1{,}2$~мм, \cite{Paxinos2019}). Через одну неделю имплантировалась GRIN-линза диаметром 1~мм (GrinTech, Германия) на 200~мкм выше CA1. Линза крепилась к черепу с помощью цианоакрилатного клея (Henkel, Германия) и зубного цемента (Stoelting, США) и покрывалась Kwik-Sil (World Precision Instruments, США). В ходе третьей операции, проведённой через неделю, над местом имплантации устанавливалась алюминиевая базовая пластина для крепления минископа, которая фиксировалась зубным цементом. Для защиты линзы устанавливался пластиковый колпачок. После финальной операции мыши возвращались в домашние клетки, и им предоставлялся трёхнедельный восстановительный период.

\subsection{Поведенческая процедура}

Через две недели после финальной операции мыши постепенно приучались к обращению и ношению минископа. В течение трёх последовательных дней каждая мышь помещалась в домашнюю клетку с подключённым, но не записывающим минископом на 5 минут за сессию.

После привыкания животные тестировались в квадратной арене открытого поля ($44 \times 44 \times 44$~см, серый пластик; Open Science, Россия) в течение четырёх последовательных дней. Каждая сессия длилась 10 минут. Арена была оборудована полупрозрачным силиконовым ковриком и четырьмя напольными объектами: белым пенопластовым конусом, флаконом для культур клеток, наполненным подстилкой, синим деревянным цилиндром и жёлтым пластиковым параллелепипедом с отверстием. На каждой стене арены была размещена уникальная визуальная метка: (1) жёлтый круг с пятью маленькими синими точками, (2) синий треугольник, (3) чёрный квадрат с широкими белыми полосами, (4) составная фигура из синих трапеций с чёрным перевёрнутым треугольником и белым кругом в центре.

В течение каждой пробы поведение животного записывалось верхней видеокамерой (Flir Chameleon3, KVR, Великобритания), одновременно с регистрацией активности CA1 минископом. Видеопоток и поток минископа синхронизировались в Bonsai (Bonsai Foundation, Великобритания) с помощью общих временных меток.

\subsection{Анализ поведения}

Отслеживание и аннотация поведения выполнялись с помощью пакета Sphynx для анализа исследовательского поведения~\cite{Plusnin1}. Для каждого видеозаписи положения ключевых точек тела оценивались покадрово с помощью нейронной сети, обученной в DeepLabCut~\cite{Mathis}. Координаты с показателем уверенности ниже 0{,}95 исключались и восстанавливались кусочно-кубической интерполяцией по соседним высокоуверенным кадрам. Полученные траектории сглаживались фильтром Савицкого--Голея третьего порядка: с окном 0{,}25~с для центра тела и основания хвоста, и 0{,}1~с для остальных точек.

Арена была сегментирована на функционально различные пространственные зоны: стеновая зона шириной 7~см по периметру, центральная зона $30 \times 30$~см, угловые зоны $7 \times 7$~см и зоны объектов шириной 2{,}5~см вокруг каждого объекта. На основе положения и движения животного вычислялись непрерывные и дискретные поведенческие переменные. К непрерывным переменным относились:
\begin{itemize}
    \item декартовы координаты центра тела;
    \item абсолютная скорость, рассчитанная по смещению центра тела и сглаженная с окном 0{,}25~с;
    \item направление тела, определяемое как угол от центра тела к центру головы;
    \item направление головы, определяемое как угол от центра головы к кончику носа.
\end{itemize}

Дискретные поведенческие акты классифицировались на основе динамики движения и пространственных переходов. Локомоторные состояния определялись по текущей скорости: быстрая локомоция ($>5$~см/с), медленная локомоция ($1{-}5$~см/с) и покой ($<1$~см/с). Замирание определялось как одновременное снижение скорости центра тела ($<1$~см/с) и скорости кончика носа ($<2$~см/с). Стойки на задних лапах (rearing) выявлялись по транзиторному сближению задних конечностей и основания хвоста. Входы в зоны и взаимодействия с объектами определялись по положению центра тела и кончика носа соответственно. Все дискретные переменные сглаживались медианным фильтром с окном 0{,}25~с для удаления шума и обеспечения минимальной длительности событий.

\subsection{Обработка данных кальциевого имиджинга}

Данные кальциевого имиджинга обрабатывались с помощью BEARMIND --- специализированного конвейера анализа, разработанного авторами (см.\ раздел~\ref{sec:driada}), построенного на основе фреймворка CaImAn~\cite{Pnevmatikakis2016}. Программное обеспечение предназначено для сквозной обработки многосессионных записей минископа.

Каждая сессия начиналась с визуального осмотра для определения границ обрезки. Предобработка включала коррекцию фона и удаление артефактов движения алгоритмом NoRMCorre~\cite{Pnevmatikakis2017}. Затем видеозаписи разбивались на пространственные маски и временные ряды активности методом CNMF-E~\cite{Zhou2018}. Гиперпараметры CNMF-E подбирались с использованием заданных критериев качества (компактность пространственной маски, SNR временного ряда и структура остатков) и верифицировались визуальным осмотром. Основные диапазоны параметров:
\begin{itemize}
    \item \texttt{gSig} (оценка размера нейрона): 3--6 пикселей (при разрешении 3~мкм/пиксель);
    \item \texttt{min-corr} (порог локальной корреляции): 0{,}85--0{,}95;
    \item \texttt{min-pnr} (отношение пикового сигнала к шуму): 3--15.
\end{itemize}

Ключевые результаты были качественно устойчивы в данном диапазоне параметров (анализ робастности); все финальные компоненты прошли ручную верификацию для исключения артефактов и дубликатов.

После декомпозиции выделенные компоненты автоматически фильтровались по порогам \texttt{rval-thr} = 0{,}95 и \texttt{min-SNR} = 3, в соответствии с рекомендациями CaImAn. Все сохранённые компоненты визуально проверялись в BEARMIND для выявления и исключения артефактов и дубликатов. Только компоненты, прошедшие ручную верификацию, классифицировались как нейроны.

Сигналы флуоресценции нормализовались как $dF/F$. Пространственные карты регистрировались между сессиями с помощью CellReg~\cite{Sheintuch2017} для продольного отслеживания индивидуальных нейронов.

\subsection{Детекция кальциевых событий}

Для извлечения дискретных кальциевых транзиентов из сигналов $dF/F$ использовался вейвлет-метод, ранее описанный в~\cite{Neugornet2021}, с незначительными модификациями. Данный метод был выбран за его устойчивость к шуму и способность обнаруживать мультимасштабные особенности в данных флуоресценции.

Сигнал $dF/F$ предварительно сглаживался свёрткой с гауссовым ядром ($\sigma = 8$), затем преобразовывался с помощью непрерывного вейвлет-преобразования (CWT), что давало двумерный массив вейвлет-коэффициентов по времени и масштабу. Значимые события проявлялись как <<гребни>> локальных максимумов на соседних масштабах.

Использовались обобщённые вейвлеты Морса~\cite{Lilly2012} $\Psi_{\beta,\gamma}(\omega) = \alpha_{\beta,\gamma}\omega^{\beta}e^{-\omega^{\gamma}}$ с параметрами $\gamma = 3$ и $\beta = 2$, соответствующими так называемым вейвлетам Эйри, которые обеспечивают минимальную площадь Гейзенберга и хорошую временнýю локализацию.

Для каждого масштаба определялись локальные максимумы вейвлет-амплитуды. Максимумы на наибольшем масштабе инициировали гребни, которые затем продлевались на меньшие масштабы по критерию временнóго совпадения. При наличии нескольких кандидатов на данном шаге сохранялся только наивысший локальный максимум в пределах каждого гребня, остальные инициировали новые гребни. Процесс продолжался до наименьшего масштаба, после чего гребни, не охватывающие все масштабы, отбрасывались.

После построения всех гребней применялся фильтр значимости по длительности гребня, пиковой амплитуде и масштабу, на котором достигался максимум. Гребни, прошедшие фильтрацию, классифицировались как кальциевые события.
